\subsection{\label{sec:qsection}qsection.txt}
%
\subsubsection{Purpose}
This file is used to specify the sections over which the fluxes are to be computed during the simulation.
A qsection is made up of all the edges associated with the same qsection name.

For each edge, the order in which the nodes are specified affects the result: when an edge is defined from its south node to its north node, a positive discharge crosses the edge from west to east.

%
\subsubsection{Format specification}
\begin{footnotesize}
\begin{tabular}{p{4cm} | p{10cm} }
\hline
Field(s) & Specification \\
\hline
\hline
$N$ & number of edges in the present file \\
\hline
$b_1 \rightharpoonup e_1 \rightharpoonup name_1$ &
\begin{itemize}
\item [$b_1$] Id of the beginning node of the 1st edge
\item [$e_1$] ID of the end node of the 1st edge in the file
\item [$name_1$] name of the 1st qsection
\item [$\rightharpoonup$] Tabulation character or space
\end{itemize} \\
\hline
$\vdots$ & (repeat) \\
\hline
$b_N \rightharpoonup e_N \rightharpoonup name_M$ & 
\begin{itemize}
\item [$b_N$] Id of the beginning node of the $N$th edge
\item [$e_N$] ID of the end node of the $N$th edge in the file
\item [$name_M$] name of the $M$th qsection
\end{itemize} \\
\hline
\end{tabular}
\end{footnotesize}

%
\subsubsection{Sample file}
In the following file, the first 5 edges are part of q section names qsection\_1, while the remaining 5 are part of q section names qsection\_2. 
\begin{verbatim}
10
1995	1741	qsection\_1
1741	1742	qsection\_1
1742	1743	qsection\_1
1743	1744	qsection\_1
1744	1745	qsection\_1
1745	1746	qsection\_2
1746	1747	qsection\_2
1747	1938	qsection\_2
1938	1939	qsection\_2
1939	1940	qsection\_2
\end{verbatim}